\chapter{Opracowanie i ocena eksperymentalna metody wypełniania żądań REST}

\section{Inżynieria zapytań}

Proces opracowania metody automatycznego mapowania wypowiedzi użytkownika w języku naturalnym na strukturalne żądania w formacie JSON, zgodne z interfejsem programistycznym aplikacji typu REST, miał charakter iteracyjny. Takie podejście jest zgodne z iteracyjnym paradygmatem rozwoju systemów generatywnej sztucznej inteligencji (\english{Generative AI}), opisywanym w literaturze jako kluczowy element poprawy jakości wnioskowania modeli LLM \cite{brown2020language,ouyang2022training}. Ewolucja zastosowanych metod obejmowała przejście od prostych strategii, umożliwiających wnioskowanie bez dostarczonych przykładów (\english{Zero-Shot}), do bardziej zaawansowanych technik, które dzięki wykorzystaniu niewielkiej liczby reprezentatywnych przykładów (\english{Few-Shot Prompting}), pozwalają precyzyjniej ukierunkować model i znacząco zwiększyć poprawność generowanych struktur \cite{brown2020language,liu2023pre}.

W początkowej fazie eksperymentów skoncentrowano się na testowaniu możliwości \definicja{dużych modeli językowych} (\akronim{LLM}, \english{Large Language Models}) bez wykorzystania dodatkowego klasyfikatora intencji. Przyjęto wówczas założenie, że model o wystarczająco dużej liczbie parametrów będzie w~stanie samodzielnie zidentyfikować cel użytkownika, dobrać odpowiedni punkt końcowy (\english{endpoint}) oraz dokonać ekstrakcji danych. W tym celu opracowano zapytanie (\english{prompt}), które zawierało opis struktury bazy danych oraz ścisłe żądanie zwrotu odpowiedzi w formacie JSON.

Pierwotna konstrukcja zapytania charakteryzowała się znaczną długością, co wynikało z konieczności zawarcia w jednym zapytaniu pełnej wiedzy o architekturze systemu. W praktyce oznaczało to bezpośrednie włączenie do promptu zarówno aktualnego stanu bazy danych dla danej hodowli, jak i wybranych fragmentów dokumentacji OpenAPI, które odzwierciedlały definicje obiektów oraz dostępnych punktów końcowych. Na wykazie \ref{wyk:prompt_verbatim} przedstawiono strukturę tego zapytania, które składało się z kilku kluczowych sekcji:

\newpage
\begin{wykaz}[ht]
    \centering
    \begin{tcolorbox}[colback=gray!5, colframe=gray!50, arc=0mm, boxrule=0.5pt]
        \begin{verbatim}
Schemat danych wejściowych (format JSON): {schema_prompt}

Twoim zadaniem jest wygenerować obiekt w formacie JSON zgodny z powyższym schematem danych wejściowych na podstawie opisu użytkownika. 

Aktualny stan bazy danych: {stan}
Użytkownik: {prompt}

Odpowiedź w formacie JSON:
        \end{verbatim}
    \end{tcolorbox}
    \caption{Struktura pierwotnego zapytania zawierająca pełny schemat bazy danych.}
    \label{wyk:prompt_verbatim}
\end{wykaz}

\begin{itemize}
    \item \textbf{Definicja schematu danych wejściowych:} sekcja \texttt{\{schema\_prompt\}} zawiera pełny opis struktur JSON dla wszystkich obsługiwanych punktów końcowych systemu. Dzięki temu model ma pełną wiedzę o formacie danych, jednak konieczne jest przetwarzanie wielu tabel i pól, często niezwiązanych z aktualnym zapytaniem.
    \item \textbf{Instrukcja zadania:} polecenie wyraźnie nakazuje wygenerowanie obiektu JSON na podstawie opisu użytkownika, zgodnie ze schematem i aktualnym stanem bazy.
    \item \textbf{Aktualny stan bazy danych:} sekcja \texttt{\{stan\}} zawiera bieżące informacje o obiektach w systemie (np. konie, kowale, weterynarze). Model może wykorzystywać te dane do wypełnienia odpowiednich pól i odwołań w JSON.
    \item \textbf{Użytkownik:} sekcja \texttt{\{prompt\}} przekazuje opis, zapytanie lub wymagania użytkownika, które mają zostać odwzorowane w formacie JSON.
    \item \textbf{Odpowiedź w formacie JSON:} model ma zwrócić wyłącznie poprawny obiekt JSON, zgodny z definicją schematu i z ograniczeniami wynikającymi z aktualnego stanu bazy oraz opisu użytkownika. Nie należy dodawać pól ani wartości spoza schematu.
\end{itemize}

Zbiorcze wyniki tej fazy testów przedstawia tabela \ref{tab:llm_comparison1}. Analiza danych wskazuje na niską skuteczność podejścia \textit{Zero-Shot}. Uzyskane wyniki pozwalają domniemywać, że przyczyną tego stanu rzeczy jest rozproszenie mechanizmu uwagi (\english{attention}) modelu na zbyt rozbudowanym opisie schematu.

Testom eksperymentalnym poddano wybrane duże modele językowe, obejmujące zarówno modele komercyjne, jak i otwartoźródłowe. Wśród nich znalazły się: 
\begin{description}
    \item[GPT-4o] (\english{OpenAI}) -- zaawansowany model komercyjny, optymalizowany pod kątem wielomodalności i wysokiej precyzji w zadaniach logicznych. \textbf{Typ:} Zamknięty (SaaS).
    
    \item[DeepSeek-V3] (\english{DeepSeek}) -- model o otwartej architekturze, charakteryzujący się wysoką wydajnością obliczeniową i zoptymalizowanym kosztem inferencji. \textbf{Typ:} Otwartoźródłowy (\english{Open-weights}).
    
    \item[Gemini 1.5 Flash] (\english{Google}) -- model komercyjny zoptymalizowany pod kątem zaawansowanego rozumowania językowego oraz przetwarzania złożonych danych wejściowych w formacie wielomodalnym. \textbf{Typ:} Zamknięty (SaaS).
\end{description}

W celu obiektywnej oceny jakości generowanych odpowiedzi zdefiniowano trzy kluczowe miary eksperymentalne:
\begin{itemize}
    \item \textbf{Skuteczność (\english{Effectiveness}):} mierzona binarnie (0 lub 1). Wartość 1 oznacza poprawną identyfikację docelowego punktu końcowego (\english{endpoint selection}), natomiast 0 wskazuje na wybór błędnego adresu URL.
    \item \textbf{Precyzja (\english{Precision}):} mierzona w znormalizowanym zakresie od 0 do 1. W badaniu przyjęto model oceny oparty na płaskiej strukturze par klucz-wartość, co wynika z braku zagnieżdżeń w testowanych schematach. Miara ta obliczana jest jako stosunek liczby poprawnie odwzorowanych elementów do sumy wszystkich unikatowych kluczy obecnych zarówno w~obiekcie wzorcowym, jak i wygenerowanym. W praktyce oznacza to, że wynik jest obniżany w~trzech przypadkach: gdy model pominie wymagany klucz, gdy przypisze mu błędną wartość lub gdy wygeneruje dodatkowe, nadmiarowe klucze spoza schematu. Wartość 1 oznacza idealną zgodność struktury i treści, natomiast 0 oznacza całkowity brak wspólnych elementów.
    \item \textbf{Efektywność (\english{Efficiency}):} wskaźnik binarny weryfikujący poprawność operacyjną wygenerowanego żądania. Wartość 1 (sukces) przypisywana jest w przypadku otrzymania kodu odpowiedzi z grupy 2xx protokołu HTTP, co potwierdza pomyślne przejście walidacji schematu oraz trwały zapis danych w bazie. Wartość 0 oznacza błąd po stronie klienta (np.~błąd struktury 4xx) lub błąd serwera (5xx), uniemożliwiający realizację żądania.
\end{itemize}

%Precyzja i efektywność liczona tylko dla poprawnie zdefiniowanych endpointów, w przeciwnym przypadku nie była liczona czy jak robimy?

Wszystkie testy eksperymentalne przeprowadzono w oparciu o zestaw przygotowanych endpointów REST, odpowiadających głównym obszarom funkcjonalnym systemu. Wśród obsługiwanych endpointów znalazły się: \texttt{/api/konie} (zarządzanie stadem), \texttt{/api/wydarzenia/rozrody} (rozród), \texttt{/api/wydarzenia/zdarzenia\_profilaktyczne} (profilaktyka), \texttt{/api/wydarzenia/podkucie} (pielęgnacja kopyt), \texttt{/api/wydarzenia/leczenia} i \texttt{/api/wydarzenia/choroby} (weterynaria i zdrowie) oraz \texttt{/api/kowale} i \texttt{/api/weterynarze} (zarządzanie kontaktami do specjalistów).

Dla każdego z tych endpointów opracowano 20 przykładowych scenariuszy użycia, które stanowiły zestaw danych testowych, a nie treningowych. Scenariusze te zostały przygotowane ręcznie przez członków zespołu badawczego i obejmowały typowe zapytania użytkowników, w tym zarówno poprawne, jak i potencjalnie problematyczne przypadki danych. Dzięki temu wyniki eksperymentów odzwierciedlają zdolność modeli do generowania poprawnych struktur JSON w podejściu Zero-Shot, bez wcześniejszego uczenia na konkretnych przykładach.

Aby zwiększyć stabilność wyników i ograniczyć wpływ losowości generatywnych modeli, każde zapytanie testowe zostało wykonane pięciokrotnie dla wszystkich badanych modeli. Ostateczne wyniki przedstawiono jako średnią z tych pięciu powtórzeń, co pozwala czytelnikowi na ocenę powtarzalności odpowiedzi modeli oraz zapewnia wiarygodność prezentowanych miar skuteczności, precyzji i efektywności. Takie podejście pozwoliło zarówno na obiektywną ocenę jakości generowanych struktur JSON, jak i na identyfikację potencjalnych problemów związanych z halucynacjami modeli czy niejednoznacznością interpretacji zapytań użytkowników.

\begin{table}[h!]
\centering
\caption{Wyniki wstępnej oceny eksperymentalnej modeli LLM (podejście Zero-Shot)}
\label{tab:llm_comparison1}
\begin{tabular}{|l|c|c|c|c|}
\hline
\textbf{Model} & \textbf{Skuteczność} & \textbf{Śr. Precyzja} & \textbf{Efektywność} \\ \hline
GPT-4o         & 0,34                 & 0,28                  & 0,41                 \\ \hline
DeepSeek-V3    & 0,45                 & 0,34                  & 0,45                 \\ \hline
Gemini 1.5 Flash & 0,41               & 0,39                  & 0,44                 \\ \hline
\end{tabular}
\end{table}

Wraz z rozwojem badań nad modelami generatywnymi istotnym elementem stało się opracowanie skutecznych technik inżynierii zapytań (\english{prompt engineering}). W literaturze \cite{liu2023pre,zhou2023comprehensive} podkreśla się, że poprawne konstruowanie zapytań może znacząco wpływać na jakość odpowiedzi LLM, szczególnie w zadaniach wymagających precyzyjnej struktury odpowiedzi lub wnioskowania na podstawie złożonych danych. Do kluczowych technik należą m.in.: stosowanie jednoznacznych instrukcji i jasnych ograniczeń formalnych, minimalizowanie niejednoznaczności poprzez kontekstowe doprecyzowanie roli modelu, przekazywanie przykładów poprawnych odpowiedzi (metody \english{Few-Shot}), stosowanie formatów szablonowych oraz technik wymuszających strukturę wyjścia (np. „Odpowiadaj WYŁĄCZNIE w poprawnym formacie JSON”). Praktyki te pozwalają ograniczyć halucynacje modeli, zwiększyć stabilność generowanych odpowiedzi oraz poprawić zgodność wyników z oczekiwaniami systemu wykonawczego.

\newpage
\begin{wykaz}[ht]
    \centering
    \begin{tcolorbox}[colback=gray!5, colframe=gray!50, arc=0mm, boxrule=0.5pt]
        \begin{verbatim}
ROLA:
Jesteś systemem przetwarzania zapytań użytkownika na poprawne struktury JSON
zgodne z interfejsem REST. Twoim zadaniem jest dokładna analiza opisu oraz
wygenerowanie wyłącznie poprawnego obiektu JSON zgodnego ze schematem danych.

ZASADY:
1. Odpowiadaj WYŁĄCZNIE w poprawnym formacie JSON.
2. Nie dodawaj żadnych pól, kluczy ani wartości, które nie występują w schemacie.
3. Jeśli użytkownik opisuje dane niezgodne ze schematem - pomiń je.
4. Jeśli użytkownik podaje wartość niepoprawną typologicznie - dopasuj ją do schematu,
ale nie zgaduj wartości niepodanych.
5. Nie generuj tekstu objaśniającego ani komentarzy.
6. Wybierz najbardziej pasujący endpoint REST na podstawie struktury danych.
7. Analizuj problem krok po kroku (nie wyświetlaj rozumowania) i dopiero potem wygeneruj JSON.
8. Jeśli dane mają odwołanie do istniejących obiektów, używaj wartości ze stanu bazy.
9. W przypadku niejednoznaczności wybierz najprostsze poprawne rozwiązanie.


SCHEMAT DANYCH (JSON):
{schema_prompt}

AKTUALNY STAN BAZY:
{stan}

UŻYTKOWNIK:
{prompt}

WYNIK:
Zwróć obiekt JSON zawierający:
{
"endpoint": "<wybrany_endpoint>",
"data": <wygenerowany_json>
}
        \end{verbatim}
    \end{tcolorbox}
    \caption{Struktura pierwotnego zapytania zawierająca pełny schemat bazy danych.}
    \label{wyk:zero-shot-prompt-engineering}
\end{wykaz}

\begin{table}[h!]
\centering
\caption{Wyniki eksperymentalne z udoskonalonym promptem Zero-Shot (wybrany endpoint)}
\label{tab:llm_comparison2}
\begin{tabular}{|l|c|c|c|}
\hline
\textbf{Model} & \textbf{Skuteczność} & \textbf{Śr. Precyzja} & \textbf{Efektywność} \\ \hline
GPT-4o & 0,62 & 0,57 & 0,65 \\ \hline
DeepSeek-V3 & 0,55 & 0,51 & 0,57 \\ \hline
Gemini 1.5 Flash & 0,59 & 0,54 & 0,60 \\ \hline
\end{tabular}
\end{table}


Analiza wyników wskazuje, że udoskonalony prompt Zero-Shot znacząco poprawił skuteczność i precyzję generowanych odpowiedzi w porównaniu z pierwotnym podejściem (Tabela \ref{tab:llm_comparison1}). Modele rzadziej generowały nieistniejące identyfikatory i lepiej odwzorowywały strukturę danych zgodnie z wybranym schematem endpointu. Efektywność, rozumiana jako poprawne wykonanie żądania w systemie REST, także uległa istotnej poprawie, co pokazuje, że zmniejszenie rozproszenia uwagi i ograniczenie promptu do konkretnego endpointu miało bezpośredni wpływ na stabilność wyników.

Wstępne wyniki pokazały, że podejście \textit{Zero-Shot} generuje liczne halucynacje, szczególnie w~zakresie identyfikatorów relacyjnych. Modele często tworzyły nieistniejące numery ID lub błędnie interpretowały polskie nazewnictwo fachowe stosowane w hodowli koni, co bezpośrednio obniżało miarę precyzji i efektywności. Przykładowo, model potrafił poprawnie wybrać endpoint (skuteczność=1), lecz z wstawiał losowe klucze obce, co skutkowało odrzuceniem żądania przez backend (efektywność wynosząca 0).

Doświadczenia z~podejściem bazowym stały się bezpośrednim impulsem do wdrożenia strategii \textit{Few-Shot Prompting}. Nowa architektura zapytania (Wykaz \ref{wyk:pierwotny_prompt}) została wzbogacona o~sekwencję par przykładów uczących, co pozwoliło modelowi na interpretację intencji użytkownika na podstawie konkretnych wzorców, a~nie tylko abstrakcyjnego opisu technicznego. Dzięki temu \definicja{duży model językowy} zyskał punkt odniesienia w~zakresie oczekiwanej struktury obiektu \textsf{JSON} oraz sposobu mapowania potocznych sformułowań na parametry funkcji \akronim{API}.

\begin{wykaz}[ht]
    \centering
    \begin{tcolorbox}[colback=gray!5, colframe=gray!50, arc=0mm, boxrule=0.5pt]
        \begin{verbatim}
ROLA:
Jesteś systemem przetwarzania zapytań użytkownika na poprawne struktury JSON
zgodne z interfejsem REST. Twoim zadaniem jest dokładna analiza opisu oraz
wygenerowanie wyłącznie poprawnego obiektu JSON zgodnego ze schematem danych.


ZASADY:
1. Odpowiadaj WYŁĄCZNIE w poprawnym formacie JSON.
2. Nie dodawaj żadnych pól, kluczy ani wartości, które nie występują w schemacie.
3. Jeśli użytkownik opisuje dane niezgodne ze schematem – pomiń je.
4. Jeśli użytkownik podaje wartość niepoprawną typologicznie – dopasuj ją do schematu,
ale nie zgaduj wartości niepodanych.
5. Nie generuj tekstu objaśniającego ani komentarzy.
6. Wybierz najbardziej pasujący endpoint REST na podstawie struktury danych.
7. Analizuj problem krok po kroku (nie wyświetlaj rozumowania) i dopiero potem wygeneruj JSON.
8. Jeśli dane mają odwołanie do istniejących obiektów, używaj wartości ze stanu bazy.
9. W przypadku niejednoznaczności wybierz najprostsze poprawne rozwiązanie.


SCHEMAT DANYCH (JSON):
{schema_prompt}


AKTUALNY STAN BAZY:
{stan}


% Pętla generująca przykłady few-shot:
for train in examples_train:
  Użytkownik: {train_prompt}
  WYNIK:
  Zwróć obiekt JSON zawierający:
  {
  "endpoint": "<wybrany_endpoint>",
  "data": <wygenerowany_json>
  }
% Koniec pętli

UŻYTKOWNIK:
{prompt}


WYNIK:
Zwróć obiekt JSON zawierający:
{
"endpoint": "<wybrany_endpoint>",
"data": <wygenerowany_json>
}
        \end{verbatim}
    \end{tcolorbox}
    \caption{Struktura zapytania zawierająca pełny schemat bazy danych i przykłady uczące.}
    \label{wyk:few-shot}
\end{wykaz}

Wykaz \ref{wyk:few-shot} przedstawia szablon zapytania. Fragmenty ujęte w nawiasy klamrowe (np. \texttt{{stan}}) oraz instrukcje sterujące (pętla \texttt{for}) stanowią elementy dynamiczne, które są automatycznie uzupełniane danymi z systemu przed wysłaniem zapytania do modelu. Struktura zoptymalizowanego zapytania w~podejściu \textit{Few-Shot} zachowała kluczowe sekcje z~poprzedniej iteracji, wzbogacając je~o~moduł demonstracyjny dla modelu. Elementy takie jak definicja schematu danych oraz ogólna instrukcja zadania pozostały identyczne jak w~przypadku metody \textit{Zero-Shot}, co pozwoliło na rzetelne porównanie wpływu samych przykładów na jakość generowanych odpowiedzi:

\begin{description}
    \item \textbf{Definicja schematu danych:} sekcja \texttt{\{schema\_prompt\}} zawiera kompletny opis struktury \textsf{JSON} dla wybranego endpointu, określający wymagane i opcjonalne pola oraz typy danych. Dzięki temu model ma pełną wiedzę o formacie odpowiedzi, jednocześnie ograniczając konieczność przetwarzania informacji niezwiązanych z bieżącym zapytaniem użytkownika.
    \item \textbf{Instrukcja zadania:} jasne polecenie nakazuje modelowi wygenerowanie poprawnego obiektu \textsf{JSON} w oparciu o opis użytkownika oraz dane ze stanu bazy, bez dodawania pól ani wartości spoza schematu. Model ma również korzystać z dostępnych danych referencyjnych w bazie przy wypełnianiu powiązań.
    \item \textbf{Przykłady uczące (Few-Shot):} sekcja iteracyjnie wstawiająca pary \texttt{\{train\_prompt\}} i \texttt{\{train\_json\}}. Te przykłady służą jako wzorzec poprawnego mapowania opisów użytkownika na struktury \textsf{JSON}, pomagając modelowi uniknąć błędów logicznych i halucynacji przy generowaniu nowych zapytań.
    \item \textbf{Aktualny stan bazy danych:} sekcja \texttt{\{stan\}} przekazuje bieżące informacje o obiektach systemu (np. konie, kowale, weterynarze), które mogą być wykorzystywane do wypełniania powiązań i wartości w generowanym \textsf{JSON}.
    \item \textbf{Użytkownik:} sekcja \texttt{\{prompt\}} zawiera aktualne zapytanie lub opis użytkownika, które mają zostać odwzorowane w formacie \textsf{JSON}.
    \item \textbf{Wynik:} model zwraca wyłącznie poprawny obiekt \textsf{JSON}, zgodny ze schematem i przykładowymi wzorcami, z uwzględnieniem aktualnego stanu bazy. Endpoint jest już wybrany przez klasyfikator, więc model generuje jedynie poprawną strukturę danych.
\end{description}

\begin{table}[h!]
\centering
\caption{Wyniki wstępnej oceny eksperymentalnej modeli LLM (podejście Few-Shot)}
\label{tab:llm_few_shot}
\begin{tabular}{|l|c|c|c|c|}
\hline
\textbf{Model} & \textbf{Skuteczność} & \textbf{Śr. Precyzja} & \textbf{Efektywność} \\ \hline
GPT-4o         & 0,74                 & 0,68                  & 0,79                 \\ \hline
DeepSeek-V3    & 0,81                 & 0,64                  & 0,65                 \\ \hline
Gemini 1.5 Flash & 0,81               & 0,84                  & 0,79                 \\ \hline
\end{tabular}
\end{table}

Zbiorcze wyniki oceny modeli w~podejściu \textit{Few-Shot} (Tabela \ref{tab:llm_few_shot}) potwierdziły znaczącą przewagę tej metody nad próbami \textit{Zero-Shot}. Na szczególną uwagę zasługują wyniki modelu \textbf{Gemini 1.5 Flash}, który wykazał się najwyższą średnią precyzją ($0,84$), zachowując przy tym wysoką skuteczność ($0,81$) oraz satysfakcjonujący wskaźnik sukcesu ($0,79$). Połączenie tych cech czyni go optymalnym silnikiem dla projektowanego \definicja{Asystenta NLP} w~kontekście stabilności struktury wyjściowej JSON.

Warto zaznaczyć, że modele \textit{Gemini 1.5 Flash} oraz \textit{DeepSeek-V3} posiadały istotną przewagę implementacyjną w stosunku do modelu \textit{GPT-4o}. Ze względu na brak dostępu do interfejsu programistycznego dla tego ostatniego, testy mogły być przeprowadzane wyłącznie w trybie interaktywnym za pośrednictwem przeglądarki. Pozostałe modele umożliwiały natomiast pełną automatyzację oraz integrację z potokiem przetwarzania danych, a także precyzyjną kontrolę parametrów generowania odpowiedzi, takich jak stopień losowości czy szczegółowość wygenerowanego tekstu. W kontekście systemów produkcyjnych działających w czasie rzeczywistym możliwość programistycznego wywołania modelu okazała się kluczowym czynnikiem przy wyborze końcowej technologii.

W toku prac wdrożeniowych nastąpiła konieczność migracji z modelu \textit{Gemini 1.5 Flash} na nowszą wersję \textit{Gemini 2.5 Flash}. Zmiana ta była podyktowana decyzją dostawcy o wyłączeniu dostępu do interfejsu programistycznego starszej wersji modelu. Migracja, choć wymuszona czynnikami zewnętrznymi, pozwoliła na utrzymanie dotychczasowych wskaźników precyzji przy jednoczesnym zapewnieniu stabilności infrastrukturalnej systemu w długofalowej perspektywie eksploatacyjnej. Dzięki tak zaprojektowanej architekturze, opartej na modelu Gemini, system osiągnął wysoką deterministyczność przy zachowaniu naturalności interfejsu użytkownika.

W praktyce podejście to generowało kilka istotnych ograniczeń inżynierskich. Ze względu na wstrzykiwanie pełnej specyfikacji bazy danych do zapytania, jego długość była znaczna, a wraz ze wzrostem liczby obsługiwanych endpointów oraz rozrostem systemu, metoda stawała się coraz mniej skalowalna.

Dodatkowo, model musiał przetwarzać rozbudowany zbiór danych wejściowych, z którego znaczna część stanowiła informacyjny „szum”. Na przykład, przy prostym poleceniu typu „dodaj nowego konia”, model analizował jednocześnie schematy związane z leczeniem, kowalami, rozrodem czy profilaktyką, co prowadziło do nadmiernego obciążenia i zwiększało ryzyko pomyłki przy wyborze właściwego endpointu lub błędnej interpretacji danych.

Dodatkowym obciążeniem było wstrzykiwanie aktualnego stanu bazy danych dla wszystkich encji jednocześnie. Dane te nie tylko konsumowały dostępne okno kontekstowe (\english{context window}), ale również prowadziły do rzadszych, lecz trudnych do zdiagnozowania błędów semantycznych, w których model próbował powiązać dane użytkownika z nieprawidłowym kontekstem pochodzącym z innej części bazy. Doświadczenia te stały się bezpośrednim impulsem do wprowadzenia dedykowanego klasyfikatora, który pozwolił na selektywne ładowanie wyłącznie tych fragmentów schematu i danych, które są krytyczne dla realizacji konkretnego żądania.